%======================================================================
\section{Track PRF --- Performance: a lighter FreeCAD}\label{sec:prf}
%======================================================================

Bruce, 2 Nov 2026: \emph{``we need to lighten the burden of FreeCAD on the system''}. The trigger was the
TL-P01 toolchanger (a Voron 2.4 350 with six tools): its working files hold about 400\,MB of shape text, and
the GUI sits on \emph{Postprocessing} for a long time while opening them.

What a CAD file holds is the same everywhere: the boundary representation (faces, edges, surfaces). How it is
written differs. FreeCAD and STEP write it as text; Parasolid (SolidWorks, NX, Onshape), ShapeManager (Fusion,
Inventor) and Granite (Creo) write binary, and store a repeated part once. The vault already does the second
(one file per part, assemblies link); PRF-01 does the first.

\begin{center}\small
\begin{tabular}{@{}lrrl@{}}
\toprule
Voron 2.4 350 (4582 objects), FreeCADCmd & Text & Binary & \\
\midrule
Open (two runs) & 2.3--2.4\,s & 1.3--1.4\,s & about 43\,\% faster \\
Save & 7.9\,s & 3.7\,s & about 53\,\% faster \\
File & 51.8\,MB & 47.5\,MB & the zip already compresses text \\
Total solid volume & 11\,980\,648\,mm$^3$ & identical & lossless \\
\bottomrule
\end{tabular}
\end{center}

Headless, the same file opens in about 2\,s, so most of the GUI's \emph{Postprocessing} is display work
(tessellation and the scene graph), not reading. Binary is the first step; PRF-02 measures the rest and PRF-03
removes it.

\begin{epic}{PRF-01}{Binary shapes by default}{M5}{\stProg}{Bruce 2 Nov}
FreeCAD already writes binary BREP when \code{BaseApp/Preferences/Document/SaveBinaryBrep} is set (stock: off;
the auto-saver already uses binary). The BruceWorx pack now sets it, the BruceWorx preference page has the
checkbox (\emph{Save shapes in binary}), and \code{bruceworx\_pack.top\_up} brings keys added to the pack later
to profiles that applied it earlier, recording their old value first so \emph{Off} still restores it. Reading
needs nothing: FreeCAD restores both forms.
\end{epic}

\begin{epic}{PRF-02}{Where the GUI time goes}{M5}{\stTodo}{Bruce 2 Nov}
Profile the GUI open of TL-P01: link restore, recompute, tessellation, Coin scene graph, tree view. Publish
each phase's share with a script that reproduces it.
\end{epic}

\begin{epic}{PRF-03}{Light display of big assemblies}{M5}{\stTodo}{Bruce 2 Nov}
From PRF-02's numbers: tessellate hidden objects only when they are shown; one tessellation per linked part
instead of one per link; open large assemblies on the preview meshes the vault already stores (PDM-13) and load
a part's exact shape when it is edited or measured. Target: TL-P01 opens at least three times faster and orbits
smoothly.
\end{epic}

\begin{epic}{PRF-04}{Vault storage migration}{M5}{\stProg}{Bruce 2 Nov}
Bruce: switch every vault file to binary, without check-outs or a new version per item. Stored files are
write-once and verified by SHA-256 (\code{objects/<ab>/<sha256>}), and drawings, BOMs and previews refer to those
hashes, so files are not edited in place. \code{bruceworx.vault\_storage} re-encodes each stored version into a
new object: a zip rewrite in which every shape Document.xml names as \code{X.brp} becomes \code{X.bin} and
everything else is copied byte for byte. Each shape is compared with the original (type, counts, volume, area,
box) and any difference stops the run before anything is recorded. One \code{migrate\_storage} journal event
maps old to new; replay moves versions, drawings' made-from, released files and BOM versions to the new hashes,
\code{FolderVault.canonical} keeps old work-copy bases current, and unchanged work copies are refreshed. Old
objects stay, so an older Vector Labs still reads the vault. On a copy of the vault: 114 of 165 stored files,
767 shapes, 26.9 to 18.5\,MB, statuses unchanged, 37 of 37 files open in FreeCAD with identical shapes.
\end{epic}
